% GENERATED FILE. Edit canonical lessons, then run npm run generate:books.
% canonical-source-hash: fnv1a64:018fab0a147b0df0
% canonical-lessons: KA-C229-hande, KA-C229-totti, KA-C229-tottilu, KA-C229-gantu, KA-C229-sarapali, KA-R229-first-pass-words-from-chapters-225-226, KA-R229-second-pass-words-from-chapters-227-229

\chapter{Large water-heating pot, Water tank, Cradle, Knot, Chain}
\label{ch:ka-large-water-heating-pot-water-tank-cradle-knot-chain}
\hlchaptermodality{229}

\begin{chapteropening}
\textbf{By the end of this chapter:} \emph{I can say a large water-heating pot, a water tank, a cradle, a knot and a chain.}

Complete the last lesson of chapter 229: I can say a large water-heating pot, a water tank, a cradle, a knot and a chain.
\end{chapteropening}

\section[\texorpdfstring{\kn{ಹಂಡೆ} (haṇḍe)}{ಹಂಡೆ (haṇḍe)}]{\kn{ಹಂಡೆ} (haṇḍe) — a large water-heating pot}

\label{lesson:KA-C229-hande}



\begin{quote}
Before the new one: say the Kannada for an opinion, then the Kannada for an idea.
\end{quote}

\subsection*{You'll want to know: \kn{ಹಂಡೆ}}
\textbf{\kn{ಹಂಡೆ}} — \emph{haṇḍe} — \textquotedblleft{}a large water-heating pot\textquotedblright{}.

\begin{grammarlens}[title={What you've built}]
One more word for this chapter.
\end{grammarlens}

\subsection*{Your turn}
\begin{itemize}
\raggedright
  \item \emph{Say it:} \emph{haṇḍe}
  \item \emph{Say it:} \emph{haṇḍe}, once more
  \item \emph{Say it:} say \emph{haṇḍe}
  \item \emph{Recall:} say the Kannada for an opinion, then the Kannada for an idea, then say \emph{haṇḍe} again
\end{itemize}

\begin{tcolorbox}[breakable,colback=teal!4,colframe=teal!35!black,title={Before you move on}]
What does \kn{ಹಂಡೆ} mean? (\textquotedblleft{}A large water-heating pot\textquotedblright{}.) Say it once more.
\end{tcolorbox}
\section[\texorpdfstring{\kn{ತೊಟ್ಟಿ} (toṭṭi)}{ತೊಟ್ಟಿ (toṭṭi)}]{\kn{ತೊಟ್ಟಿ} (toṭṭi) — a water tank, a trough}

\label{lesson:KA-C229-totti}



\begin{quote}
Before the new one: say the Kannada for an idea, then the Kannada for a large water-heating pot.
\end{quote}

\subsection*{You'll want to know: \kn{ತೊಟ್ಟಿ}}
\textbf{\kn{ತೊಟ್ಟಿ}} — \emph{toṭṭi} — \textquotedblleft{}a water tank, a trough\textquotedblright{}.

\begin{grammarlens}[title={What you've built}]
One more word for this chapter.
\end{grammarlens}

\subsection*{Your turn}
\begin{itemize}
\raggedright
  \item \emph{Say it:} \emph{toṭṭi}
  \item \emph{Say it:} \emph{toṭṭi}, once more
  \item \emph{Say it:} say \emph{toṭṭi}
  \item \emph{Recall:} say the Kannada for an idea, then the Kannada for a large water-heating pot, then say \emph{toṭṭi} again
\end{itemize}

\begin{tcolorbox}[breakable,colback=teal!4,colframe=teal!35!black,title={Before you move on}]
What does \kn{ತೊಟ್ಟಿ} mean? (\textquotedblleft{}A water tank, a trough\textquotedblright{}.) Say it once more.
\end{tcolorbox}
\section[\texorpdfstring{\kn{ತೊಟ್ಟಿಲು} (toṭṭilu)}{ತೊಟ್ಟಿಲು (toṭṭilu)}]{\kn{ತೊಟ್ಟಿಲು} (toṭṭilu) — a cradle}

\label{lesson:KA-C229-tottilu}



\begin{quote}
Before the new one: say the Kannada for a large water-heating pot, then the Kannada for a water tank.
\end{quote}

\subsection*{You'll want to know: \kn{ತೊಟ್ಟಿಲು}}
\textbf{\kn{ತೊಟ್ಟಿಲು}} — \emph{toṭṭilu} — \textquotedblleft{}a cradle\textquotedblright{}.

\begin{grammarlens}[title={What you've built}]
One more word for this chapter.
\end{grammarlens}

\subsection*{Your turn}
\begin{itemize}
\raggedright
  \item \emph{Say it:} \emph{toṭṭilu}
  \item \emph{Say it:} \emph{toṭṭilu}, once more
  \item \emph{Say it:} say \emph{toṭṭilu}
  \item \emph{Recall:} say the Kannada for a large water-heating pot, then the Kannada for a water tank, then say \emph{toṭṭilu} again
\end{itemize}

\begin{tcolorbox}[breakable,colback=teal!4,colframe=teal!35!black,title={Before you move on}]
What does \kn{ತೊಟ್ಟಿಲು} mean? (\textquotedblleft{}A cradle\textquotedblright{}.) Say it once more.
\end{tcolorbox}
\section[\texorpdfstring{\kn{ಗಂಟು} (gaṇṭu)}{ಗಂಟು (gaṇṭu)}]{\kn{ಗಂಟು} (gaṇṭu) — a knot; a bundle}

\label{lesson:KA-C229-gantu}



\begin{quote}
Before the new one: say the Kannada for a water tank, then the Kannada for a cradle.
\end{quote}

\subsection*{You'll want to know: \kn{ಗಂಟು}}
\textbf{\kn{ಗಂಟು}} — \emph{gaṇṭu} — \textquotedblleft{}a knot; a bundle\textquotedblright{}.

\begin{grammarlens}[title={What you've built}]
One more word for this chapter.
\end{grammarlens}

\subsection*{Your turn}
\begin{itemize}
\raggedright
  \item \emph{Say it:} \emph{gaṇṭu}
  \item \emph{Say it:} \emph{gaṇṭu}, once more
  \item \emph{Say it:} say \emph{gaṇṭu}
  \item \emph{Recall:} say the Kannada for a water tank, then the Kannada for a cradle, then say \emph{gaṇṭu} again
\end{itemize}

\begin{tcolorbox}[breakable,colback=teal!4,colframe=teal!35!black,title={Before you move on}]
What does \kn{ಗಂಟು} mean? (\textquotedblleft{}A knot; a bundle\textquotedblright{}.) Say it once more.
\end{tcolorbox}
\section[\texorpdfstring{\kn{ಸರಪಳಿ} (sarapaḷi)}{ಸರಪಳಿ (sarapaḷi)}]{\kn{ಸರಪಳಿ} (sarapaḷi) — a chain}

\label{lesson:KA-C229-sarapali}



\begin{quote}
Before the new one: say the Kannada for a cradle, then the Kannada for a knot.
\end{quote}

\subsection*{You'll want to know: \kn{ಸರಪಳಿ}}
\textbf{\kn{ಸರಪಳಿ}} — \emph{sarapaḷi} — \textquotedblleft{}a chain\textquotedblright{}.

\begin{grammarlens}[title={What you've built}]
One more word for this chapter.
\end{grammarlens}

\subsection*{Your turn}
\begin{itemize}
\raggedright
  \item \emph{Say it:} \emph{sarapaḷi}
  \item \emph{Say it:} \emph{sarapaḷi}, once more
  \item \emph{Say it:} say \emph{sarapaḷi}
  \item \emph{Recall:} say the Kannada for a cradle, then the Kannada for a knot, then say \emph{sarapaḷi} again
\end{itemize}

\begin{tcolorbox}[breakable,colback=teal!4,colframe=teal!35!black,title={Before you move on}]
What does \kn{ಸರಪಳಿ} mean? (\textquotedblleft{}A chain\textquotedblright{}.) Say it once more.
\end{tcolorbox}
\section[(dialogue)]{Words from chapters 225-226 — a first pass}

\label{lesson:KA-R229-first-pass-words-from-chapters-225-226}



\begin{quote}
Say the words for a knot and a chain.
\end{quote}

\subsection*{Your turn}
Say these aloud:

\begin{itemize}
\raggedright
  \item neem, jasmine, a sunflower, sugarcane, a firecracker
  \item art, aunt, daughter-in-law, son-in-law, sister-in-law
\end{itemize}

\begin{tcolorbox}[breakable,colback=teal!4,colframe=teal!35!black,title={Before you move on}]
Neem? (\textbf{\kn{ಬೇವು}}, \emph{bēvu}.) Jasmine? (\textbf{\kn{ಮಲ್ಲಿಗೆ}}, \emph{mallige}.) A sunflower? (\textbf{\kn{ಸೂರ್ಯಕಾಂತಿ}}, \emph{sūryakānti}.) Sugarcane? (\textbf{\kn{ಕಬ್ಬು}}, \emph{kabbu}.) A firecracker? (\textbf{\kn{ಪಟಾಕಿ}}, \emph{paṭāki}.) Art? (\textbf{\kn{ಕಲೆ}}, \emph{kale}.) Aunt? (\textbf{\kn{ಅತ್ತೆ}}, \emph{atte}.) Daughter-in-law? (\textbf{\kn{ಸೊಸೆ}}, \emph{sose}.) Son-in-law? (\textbf{\kn{ಅಳಿಯ}}, \emph{aḷiya}.) Sister-in-law? (\textbf{\kn{ಅತ್ತಿಗೆ}}, \emph{attige}.)
\end{tcolorbox}
\section[(dialogue)]{Words from chapters 227-229 — a second pass}

\label{lesson:KA-R229-second-pass-words-from-chapters-227-229}



\begin{quote}
Say the words for a saw and a ladder.
\end{quote}

\subsection*{Your turn}
Say these aloud:

\begin{itemize}
\raggedright
  \item a saw, a ladder, glue, a torch, a machine
  \item football, chess, a hobby, an opinion, an idea
  \item a large water-heating pot, a water tank, a cradle, a knot, a chain
\end{itemize}

\begin{tcolorbox}[breakable,colback=teal!4,colframe=teal!35!black,title={Before you move on}]
A saw? (\textbf{\kn{ಗರಗಸ}}, \emph{garagasa}.) A ladder? (\textbf{\kn{ಏಣಿ}}, \emph{ēṇi}.) Glue? (\textbf{\kn{ಅಂಟು}}, \emph{aṇṭu}.) A torch? (\textbf{\kn{ಟಾರ್ಚ್}}, \emph{ṭārc}.) A machine? (\textbf{\kn{ಯಂತ್ರ}}, \emph{yantra}.) Football? (\textbf{\kn{ಫುಟ್ಬಾಲ್}}, \emph{phuṭbāl}.) Chess? (\textbf{\kn{ಚದುರಂಗ}}, \emph{caduraṅga}.) A hobby? (\textbf{\kn{ಹವ್ಯಾಸ}}, \emph{havyāsa}.) An opinion? (\textbf{\kn{ಅಭಿಪ್ರಾಯ}}, \emph{abhiprāya}.) An idea? (\textbf{\kn{ವಿಚಾರ}}, \emph{vicāra}.) A large water-heating pot? (\textbf{\kn{ಹಂಡೆ}}, \emph{haṇḍe}.) A water tank? (\textbf{\kn{ತೊಟ್ಟಿ}}, \emph{toṭṭi}.) A cradle? (\textbf{\kn{ತೊಟ್ಟಿಲು}}, \emph{toṭṭilu}.) A knot? (\textbf{\kn{ಗಂಟು}}, \emph{gaṇṭu}.) A chain? (\textbf{\kn{ಸರಪಳಿ}}, \emph{sarapaḷi}.)
\end{tcolorbox}
